\documentclass[a4paper,12pt]{article}

\usepackage[utf8]{inputenc}
\usepackage[english]{babel}
\usepackage{amsmath}
\usepackage{graphicx}
\usepackage{csvsimple}
\usepackage{booktabs}
\usepackage{caption}

\title{Hotspot Analysis}
\author{Andrea Cattarinich S5137057}
\date{\today}

\begin{document}

\maketitle

\section{Hotspot Analysis}

The goal of this analysis is to identify the parts of the program that consume the most execution time.  
The measurements were performed using two methods: 
\begin{itemize}
    \item \textbf{Manually}: using \texttt{omp\_get\_wtime()} 
    \item \textbf{Profiler}: using \texttt{Intel VTune}
\end{itemize}
Either two methods have been used with with [$N = 10000$] on [WSL2-workstation] and produced [similar results].

\begin{center}
    \begin{minipage}{0.48\textwidth}
    \centering
    \captionof{table}{Manually}
    \resizebox{\textwidth}{!}{%
    \begin{tabular}{lrr}
        \toprule
        Function & Time (s) & Percentage (\%) \\
        \midrule
        \csvreader[late after line=\\]{hotspots_10000_manually.csv}{}
        {\csvcoli & \csvcolii & \csvcoliii} 
        \bottomrule
    \end{tabular}%
    }
    \end{minipage}%
    \hfill
    \begin{minipage}{0.48\textwidth}
    \centering
    \captionof{table}{Intel VTune Profiler}
    \resizebox{\textwidth}{!}{%
    \begin{tabular}{lrr}
        \toprule
        Function & Time (s) & Percentage (\%) \\
        \midrule
        \csvreader[late after line=\\]{hotspots_10000_intel_profiler.csv}{}
        {\csvcoli & \csvcolii & \csvcoliii} 
        \bottomrule
    \end{tabular}%
    }
    \end{minipage}
\end{center}

\noindent
% TODO:
Aggiungere i NUMBERS e le linee xx-yy del codice



\subsection{Conclusion}
The main hotspot is the DFT() function and its inverse IDFT(), which together account for over 99\% of the total execution time. All other functions have negligible impact. This indicates that \textbf{parallelization and vectorization efforts should focus on the double loop inside DFT()}.

\end{document}